% Generated draft: requires mathematical and source review.
\documentclass[11pt]{article}
\usepackage[margin=25mm]{geometry}
\usepackage{fontspec}
\setmainfont{Latin Modern Roman}
\usepackage{amsmath,amssymb,mathtools,unicode-math}
\setmathfont{Latin Modern Math}
\usepackage{enumitem,needspace,xurl,hyperref}
\setlength{\parindent}{0pt}
\setlength{\parskip}{5pt}
\setlength{\emergencystretch}{3em}
\newcommand{\N}{\mathbb{N}}
\newcommand{\Z}{\mathbb{Z}}
\newcommand{\Q}{\mathbb{Q}}
\newcommand{\R}{\mathbb{R}}
\newcommand{\C}{\mathbb{C}}
\newcommand{\F}{\mathbb{F}}
\newcommand{\A}{\mathbb{A}}
\newcommand{\PP}{\mathbb{P}}
\newcommand{\E}{\mathbb{E}}
\newcommand{\eps}{\varepsilon}
\newcommand{\dd}{\,\mathrm{d}}
\begin{document}
\section*{UnsolvedMath ID: 1877}
\textbf{Source problem code:} catalogue entry
\subsection{Definitions and mathematical statement}
For integers $P,Q$, define the Lucas sequence of the second kind by
\[
V_0(P,Q)=2,\qquad V_1(P,Q)=P,\qquad
V_{j+2}(P,Q)=P V_{j+1}(P,Q)-Q V_j(P,Q)\quad(j\geq0).
\]
For a specified integer parameter $m$, the congruence in the supplied background is
\[
V_n(m,-1)\equiv m\pmod n.
\]
The background defines a Fibonacci pseudoprime of the $m$th kind as an \emph{odd composite} $n$ satisfying this congruence. Taken with that definition, its question about an even Fibonacci pseudoprime is internally incompatible: the same $n$ would be required to be both odd and even. This is a formulation defect rather than a number-theoretic existence problem.

If the intended definition is amended to permit even composite integers, a parameter-specific candidate question is
\[
\exists n\in\Z_{\geq4}:\quad 2\mid n,\quad n\text{ composite},\quad
V_n(m,-1)\equiv m\pmod n,
\]
where $m$ must first be fixed and any additional pseudoprime hypotheses must be supplied. Asking this for a particular $m$, for every $m$ in a stated range, or for at least one $m$ gives different targets. The stored statement provides no parameter or range and no explicit identification between the historically named ``first kind'' convention and the displayed parameterized definition. No such identification is silently imposed here.

The background attributes to Somer the earlier constraint that an even Fibonacci pseudoprime, if it exists, exceeds $28\times10^{12}$, and cites Richard Guy's Section A12. The supplied literature assessment, checked 2026-08-17, reports that Di Porto proved nonexistence of even Fibonacci pseudoprimes of the first kind and that the same primary source records even examples for other parameter choices. It therefore assesses the historically standard first-kind reading as solved negatively but labels the unqualified imported record partially solved because the convention remains unspecified. No concrete generalized even example or its exact parameter conditions is supplied. The stated remaining task is to select the Lucas parameters and remove the oddness restriction before applying a theorem or example. The primary registry status remains ``open''.
\subsection{Short English statement}
The source asks whether an even Fibonacci pseudoprime exists but defines its pseudoprimes to be odd. After that defect is repaired, the answer depends on the Lucas parameters: the supplied assessment reports nonexistence for the historical first-kind convention and even examples for other choices.
\subsection{Sources}
\begin{itemize}
\item UnsolvedMath local registry, permanent ID 1877, source code catalogue entry, ``Even Fibonacci Pseudoprimes''. The background attributes the problem to Richard Guy, \emph{Unsolved Problems in Number Theory}, Section A12, without edition or page details. Dataset attribution: Ulam AI and the UnsolvedMath Contributors; curation and metadata under CC BY 4.0.
\item Adina Di Porto, \emph{Nonexistence of Even Fibonacci Pseudoprimes of the 1st Kind}, Fibonacci Quarterly 31 (1993), 166--172, \url{https://www.fq.math.ca/Scanned/31-2/diporto.pdf}. The supplied triage cites first-kind nonexistence and discussion of parameter-dependent even examples.
\item Lawrence Somer, \emph{On Even Fibonacci Pseudoprimes}, Applications of Fibonacci Numbers 4 (1991), 353--366, \url{https://doi.org/10.1007/978-94-011-3586-3_31}. The supplied assessment associates this with the earlier first-kind constraints mentioned in the background.
\item The local literature triage records checks on 2026-08-17; no cited source was independently visited for this draft.
\end{itemize}
\end{document}
